\documentclass{article}
\usepackage{amsmath,amssymb}
\begin{document}
\section*{Derivation 3: Spin-1/2 as Temporal-Orientation Holonomy}

The temporal field $\tilde{g}_{\mu\nu}$ carries an orientation bundle with structure group $O(1,3)$. The local matter-clock congruence defines a preferred temporal direction at each point.

The conformal sector gives zero holonomy: $\oint_C \Sigma_\mu \, dx^\mu = \oint_C d(\ln A) = 0$ for any contractible loop where $A(\phi)$ is smooth and single-valued, in agreement with Paper 0 (\S3.3).

The holonomy resides in the postulated orientation sector. Around a closed loop $C$ encircling the defect core:
\[
\Delta \phi = \oint_C \vartheta_\mu \, dx^\mu = n \cdot 4\pi, \quad n \in \mathbb{Z}
\]

The matter-clock phase couples with half-integer weight, so the minimal winding $n = \pm 1$ preserves single-valuedness of $\Psi = e^{iS/\hbar}$. The two orientations correspond to:
\begin{itemize}
\item \textbf{Spin up:} Orientation phase winds with local shear ($n = +1$)
\item \textbf{Spin down:} Orientation phase winds against local shear ($n = -1$)
\end{itemize}

$SU(2)$ is proposed as the holonomy group of the orientation bundle. The Pauli matrices are generators of infinitesimal rotations in this bundle, not in physical space. The $U(1)$ winding is the abelian core of this proposal; the non-abelian completion is developed in TEP-SPIN (Paper 24).

\textbf{Conclusion:} Spin-$\tfrac{1}{2}$ is temporal-orientation holonomy, not intrinsic angular momentum.

\end{document}
